\section*{Chapter \ref{ch:m2:market-making-games} --- Market-Making Games}

\begin{solution}{exo:m2:market-making-games:1}
Around the expected value, 3.5: for example ``3 at 4, 10 up''. The value is known, so no
one can be better informed; the width is a small profit for the risk of the roll and a
width at which the other side still wants to trade.
\end{solution}

\begin{solution}{exo:m2:market-making-games:2}
Raise the market: the buyer may know better, and at least thought 400 was cheap. Move the
centre up and perhaps widen, for example to 300 at 500, and think about why she bought:
if she might have counted the windows, move a lot.
\end{solution}

\begin{solution}{exo:m2:market-making-games:3}
An informed trader passes when the sum is between 30 and 40, and an uninformed one passes
with the same probability whatever the sum; the prior is symmetric about 35 and so is
the set of sums that produce a pass, so the mean does not move. Later, when quotes are
off-centre, a pass does move it.
\end{solution}

\begin{solution}{exo:m2:market-making-games:4}
A width equal to the patience, 8 points: $4e^{-1} = 1.47$ a round, 29.4 over 20 rounds.
\end{solution}

\begin{solution}{exo:m2:market-making-games:5}
0.31. A quarter of the sums are above 40, so informed buyers are 5\% of arrivals
($0.2 \times 0.25$); uninformed traders trade only 29\% of the time at this width and buy
half of that, 11.5\% of arrivals ($0.8 \times 0.29 / 2$). The width filters the uninformed
more than the informed, so a buy is informed more often than the table's 20\%.
\end{solution}

\begin{solution}{exo:m2:market-making-games:6}
The market maker's quote is traded against precisely when someone thinks it is wrong in
their favour; conditional on being filled, the value is worse than the maker thought.
Quoting the unconditional estimate means losing on average to those who choose to trade.
\end{solution}

\begin{solution}{exo:m2:market-making-games:7}
12.5 points, earning 13.8 a game: more informed traders push the best width up and the
profit down (11.0 points and 20.6 with 10\% informed).
\end{solution}

\begin{solution}{exo:m2:market-making-games:8}
More trades earn more spread from the uninformed, but a tight market is also traded
whenever the informed know it is wrong, and those trades cost more than the spread
earns. Below the best width, profit falls; at zero width the maker loses 22 a game in the
chapter's table.
\end{solution}

\begin{solution}{pb:m2:market-making-games:1}
\textbf{1.} Mean 35, standard deviation 8.03, from 6 to 64.
\textbf{2.} A deck has four cards of each value: five aces or five kings are impossible.
\textbf{3.} $e^{-10/8} = 0.29$.
\textbf{4.} 30 at 40.
\textbf{5.} Half the width when it trades: $0.29 \times 5 = 1.43$ per round from an
arriving uninformed trader.
\textbf{6.} 38.13 after a buy; 31.87 after a sale.
\textbf{7.} On the first round the quotes are symmetric around a symmetric prior; later
the centre has moved, and a pass says the sum is probably inside the current quotes.
\textbf{8.} It rises from 35 to about 39 after early buys, then falls towards 31 as sales
come in; the early buys were uninformed noise, or informed buys against a market still
centred too low, and the maker cannot tell which.
\textbf{9.} 16.4 a game at 11.5 points against 12.0 at 14 points without updating.
\textbf{10.} Without learning its centre stays wrong, so it must keep more distance from
the value to lose less to the informed.
\textbf{11.} 11.5 points, 16.4 a game.
\textbf{12.} 14.1 at 8 points and 15.1 at 15.
\textbf{13.} 8 points, about 29 a game.
\textbf{14.} 11.0 points with 10\% informed (20.6 a game), 12.5 with 30\% (13.8).
\textbf{15.} $e^{-11.5/8} = 0.24$.
\textbf{16.} From the mark-outs of its fills: how far the value moves against it after
trades, by counterparty and size; persistent adverse moves mean informed flow.
\textbf{17.} Raise the centre, since the buyers may know the value is higher, and consider
skewing: a higher offer to sell less while short, a bid closer to buy back.
\textbf{18.} That the centre is an estimate to be revised with every trade, the width a
price for adverse selection set against how much the uninformed will pay, and that
learning from flow is worth as much as a wider spread.
\textbf{19.} Named result: \emph{the card-sum market}: against a table with a fifth of the
traders informed, a width of 11.5 points maximises the expected P\&L, 16.4 a game.
\textbf{20.} It pays for the trades made with people who know more.
\end{solution}

\begin{solution}{iq:m2:market-making-games:1}
The sum has mean 7, from 2 to 12; with no one better informed, quote around 7, for
example ``6.5 at 7.5'' for size, wider if the size is large relative to the risk one can
take, and think about who would want to trade it.
\iqlookfor{the centre, a reasoned width, and who is on the other side.}
\end{solution}

\begin{solution}{iq:m2:market-making-games:2}
Ask why: the buyer may know more. Update the estimate upwards by an amount that depends
on how likely the buyer is informed; move the market up, perhaps widen, and consider the
inventory I now hold short.
\iqlookfor{conditional expectation and inventory.}
\end{solution}

\begin{solution}{iq:m2:market-making-games:3}
Glosten and Milgrom: a risk-neutral maker facing informed and uninformed traders sets the
offer at the expected value given a buy and the bid at the expected value given a sale;
since buys are more likely when the value is high, the offer is above the bid even with
zero costs and zero expected profit.
\iqlookfor{conditional expectations on trade direction.}
\end{solution}

\begin{solution}{iq:m2:market-making-games:4}
Move the market up, both because the sales may have been to informed buyers and because
the short position is risk: raise the bid more than the offer to attract sellers and
discourage buyers, until the inventory is back within limits.
\iqlookfor{information and inventory both push the same way.}
\end{solution}

\begin{solution}{iq:m2:market-making-games:5}
From the relation between trades and subsequent price moves: estimate how much prices
move after buys and sales (the permanent impact), or fit a Glosten--Milgrom or PIN-type
model to the sequence of buys and sales and the spread; mark-outs by counterparty give a
direct measure.
\iqlookfor{price impact and mark-outs as evidence of information.}
\end{solution}

\begin{solution}{iq:m2:market-making-games:6}
A server holding each game's hidden state and order book, with a message protocol for
quotes, trades and cancellations; bots as clients with the same interface as humans;
time-stamped logs of everything for replay and scoring; rooms and scenarios defined as
data; fairness checks on latency between players; post-game analytics of P\&L, width and
adverse selection.
\iqlookfor{authoritative state, same interface for bots and humans, replay.}
\end{solution}
